\documentclass[11pt,a4paper]{article}
\usepackage[margin=27mm]{geometry}
\usepackage{amsmath,amssymb,amsthm,booktabs,hyperref}
\newtheorem{theorem}{Theorem}
\newtheorem{lemma}[theorem]{Lemma}
\newtheorem{corollary}[theorem]{Corollary}
\newcommand{\F}{\mathbb F}
\newcommand{\Sz}{\operatorname{Sz}}
\newcommand{\gen}[1]{\langle #1\rangle}
\title{Compatible tori and a missing Suzuki geometry}
\author{Human Compression project}\date{}
\begin{document}\maketitle
\begin{abstract}
A common binary matrix algebra places two nonsplit tori in compatible
subfield groups. The resulting rank-three geometry has maximal
parabolics $\Sz(8)$, $\Sz(32)$ and $C_{533}:C_4$, and is absent from
the stated exhaustive classification. We rework the previously reviewed
explicit example, replacing the finite intersection census by a
coprime-order argument. A separate corollary of the scalar-subspace
construction contradicts the same source's rank-two uniqueness assertion.
Neither argument gives a corrected complete classification.
\end{abstract}

\section{Statement and retained inputs}
For subgroups $G_i$ of $G$, the elements of the coset pregeometry are
the labelled left cosets $gG_i$; incidence is nonempty intersection.
Write $G_J=\bigcap_{j\in J}G_j$ and $G_\varnothing=G$.
The definition on printed page 46 of \cite{source} calls this system
residually weakly primitive if for every proper subset $J$ of the type
set some unused $i$ makes $G_{J\cup\{i\}}$ a maximal proper subgroup
of $G_J$. We use exactly this definition.

\begin{theorem}\label{main}
In $G=\Sz(32768)$ there are standard subfield subgroups
$S=\Sz(8)$, $T=\Sz(32)$ and a subgroup $K=C_{533}:C_4$ with
\[
 S\cap T=C_5:C_4,\quad S\cap K=C_{13}:C_4,\quad
 T\cap K=C_{41}:C_4,\quad S\cap T\cap K=C_4.
                                                        \tag{1}
\]
Their rank-three coset geometry is residually weakly primitive, thick,
flag transitive and residually connected. Its intersection lattice is
absent from the applicable cases of Theorem 5.5 of \cite{source}.
\end{theorem}

The standard faithful Suzuki matrices are denoted by $U(a,b)$,
$h(t)$ and $w$, where $w$ is the antidiagonal identity matrix. Their
formulas, Bruhat forms and exact coefficient-field intersections are
given in \cite[Lemma 1]{subfields}. In particular, for standard
embeddings, $\Sz(2^a)\cap\Sz(2^b)=\Sz(2^{\gcd(a,b)})$,
including the degree-one group $\Sz(2)=C_5:C_4$.
We retain the classical Suzuki order formula
\[
                       |\Sz(q)|=q^2(q^2+1)(q-1),         \tag{2}
\]
maximality of simple subfield groups of prime relative field degree,
and maximality of the full normalizers of nonsplit cyclic tori of
orders $q\pm\sqrt{2q}+1$. Such a normalizer has order four times the
torus order. These structural facts are among the inputs used by
\cite[Section 3]{source}; completeness of its geometry lists is not
used here. Our common proof \cite[Lemma 1]{subfields}, rather than
a new subfield-intersection calculation, is an explicit prerequisite.

\section{One algebra, two compatible tori}
All displayed integer entries in this section encode binary polynomial
coefficients in
\[
                       F=\F_2[z]/(z^{15}+z+1).
\]
Thus, for example, $5682$ denotes
$\sum_{i=0}^{14}\epsilon_i z^i$ with binary digits $\epsilon_i$ of
$5682$; it is not an integer reduced modulo $32768$.
The polynomial is irreducible. The supplied arithmetic control also
checks this directly: successive multiplication by $z$ visits all
$32767$ nonzero polynomial classes before returning to one.

Put
\[
 M=\begin{pmatrix}1&1&1&0\\0&0&1&1\\0&1&1&1\\1&0&0&0\end{pmatrix},
 \qquad
 r=\begin{pmatrix}1&0&0&0\\1&1&0&0\\1&1&1&0\\1&0&1&1\end{pmatrix}
   =U(1,0).
\]
Direct multiplication gives
\[
 M^4+M+I=0,\quad |M|=15,\quad |r|=4,
 \qquad r^{-1}Mr=M^2.                                   \tag{3}
\]
The polynomial $X^4+X+1$ has no binary root and is not divisible by
$X^2+X+1$, the only irreducible quadratic over $\F_2$.
Consequently $\F_2[M]\cong\F_{16}$. Its scalar extension to $F$
is a field because $\gcd(4,15)=1$.

Take $\lambda_8=5682$ and $\lambda_{32}=624$ in the stated encoding,
and set
\[
 a=(M+\lambda_8 I)^{315},\qquad
 b=(M+\lambda_{32}I)^{25575}.                            \tag{4}
\]
Here $315=(8^4-1)/13$ and $25575=(32^4-1)/41$.
The following table gives their full open Bruhat parameters:
\[
\begin{array}{c|rrrrr|r}
 g&u&v&t&x&y&|g|\\\hline
 a&893&0&5682&893&893&13\\
 b&5992&4671&1131&5992&1831&41
\end{array}
\qquad g=U(u,v)h(t)wU(x,y).                              \tag{5}
\]
Equations (3)--(5), the scalar identities
$\lambda_8^8=\lambda_8$, $\lambda_{32}^{32}=\lambda_{32}$,
and the indicated orders are finite polynomial identities in $F$.
They are checked by the short witness program supplied with this note,
using repeated squaring and the displayed modulus. Each entry of the
first parameter row is fixed by eighth powering; each entry of the
second is fixed by thirty-second powering. Thus $a\in S$ and $b\in T$.
The table is also a direct membership certificate, independently of
any enumeration of a Suzuki group.

Since $a$ and $b$ are polynomials in $M$, they commute. Their distinct
prime orders give $|ab|=13\cdot41=533$ and
$\gen{ab}=\gen{a,b}$. With $c=ab$, the remaining matrix identity is
\[
                     r^{-1}cr=c^{278},\qquad
                     278^2\equiv-1\pmod{533}.            \tag{6}
\]
It follows that $K=\gen{c,r}=C_{533}:C_4$: the two cyclic factors
have coprime orders, and (6) is an automorphism of order four.
Every element has a unique form $a^i b^j r^k$, with
$0\le i<13$, $0\le j<41$ and $0\le k<4$.
In particular $|K|=2132$. This $K$ is not the full normalizer of
$C_{533}$ in the ambient group; no such maximality is asserted.

\begin{lemma}\label{meets}
The intersections are exactly those in (1).
\end{lemma}
\begin{proof}
The standard-subfield result gives $S\cap T=\Sz(2)$ of order twenty.
Both $a$ and $r$ belong to $S$. If $a^i b^j r^k\in S$, multiplication
on the left by $a^{-i}$ and on the right by $r^{-k}$ gives $b^j\in S$.
But $41$ does not divide $|S|=29120$. Since $b$ has order $41$,
we must have $j=0$. Hence $K\cap S=\gen{a,r}$, of order $52$.
Similarly $13$ does not divide $|T|=32537600$, so
$K\cap T=\gen{b,r}$, of order $164$.
Their common intersection contains $\gen{r}$ and has order dividing
$\gcd(52,164)=4$, proving the final equality.
The groups of orders $52$ and $164$ are the full nonsplit torus
normalizers in $S$ and $T$: (6) normalizes their cyclic factors, and
the retained normalizer order is respectively $4\cdot13$ and
$4\cdot41$.
\end{proof}

\section{The geometry and the source contradiction}
We verify all seven instances of residual weak primitivity. The
subgroup $S$ is maximal in $G$ because the relative field degree is
$15/3=5$. Inside $S$ and $T$, respectively, choose the maximal torus
normalizers $S\cap K$ and $T\cap K$. Inside $K$, choose $K\cap S$,
of prime index $41$. At each pairwise intersection the final subgroup
$\gen{r}$ has prime index, respectively $5$, $13$ and $41$.
These are all proper type subsets. All eight intersections are distinct,
by their orders. This verifies the exact subgroup condition of the source.

For clarity, the geometry properties can be proved without its general
rank-three assertion. We use the following elementary observation.
If $A,B,C$ are subgroups and
\[
                         A=(A\cap B)(A\cap C),           \tag{7}
\]
every triple flag of their left-coset system has a common point.
Translate the $A$-coset to $A$, and write the other two cosets as
$bB,cC$ with $b,c\in A$. Factor $b^{-1}c=xy$ with
$x\in A\cap B$, $y\in A\cap C$. Then $bx=cy^{-1}$ lies in
all three cosets. Every flag of fewer types also extends to a chamber:
an incident pair has a common point, and that point supplies the missing
coset. Thus (7) gives a geometry and transitivity on every type of flag.

In our example the product set $(K\cap S)(K\cap T)$ has cardinality
\[
                         \frac{52\cdot164}{4}=2132=|K|.
\]
It is contained in $K$, so equals $K$. No assumption that an arbitrary
subgroup product is a subgroup is needed. Apply (7) with $A=K$.
The three panel sizes are $5$, $13$ and $41$, so the geometry is thick.

For residual connectivity, the two residue subgroups in $S$ generate
$S$: one is its maximal torus normalizer of order $52$, while the
other, $S\cap T$, contains an element of order five and is not contained
in that normalizer. The same argument, using order $164$, works in $T$.
The product formula gives generation in $K$. Finally $\gen{S,T}=G$,
because $S$ is maximal and $31$ divides $|T|$ but not $|S|$.
The usual coset-graph argument now proves connectedness of the whole
incidence graph and of every rank-two residue: a component containing
the base flag consists exactly of the cosets reached by words in its
incident stabilizer subgroups. Those subgroups generate the relevant
group, as just shown. Flag transitivity transports this statement to
all residues.

It remains to compare the actual source conclusion. Here the odd field
degree is $15=3\cdot5$, whose prime exponents sum to two. Moreover
the subfield $S=\Sz(8)$ has $8-1$ prime. The applicable portion of
\cite[Theorem 5.5, printed pp.58--59]{source} is therefore its first
five shapes. Only shape 5 has two Suzuki maximal parabolics and one
odd cyclic group extended by $C_4$. It prescribes the intersection of
the two Suzuki parabolics as $C_4\times C_2$, of order eight; our
intersection has order twenty. A permutation of types cannot change
this discrepancy, since $K$ is soluble and $S,T$ are not. Shape 8
belongs to the separate branch where the sum of prime exponents
exceeds two, and cannot be used in this branch. This proves
Theorem~\ref{main}.

The source's closing qualification on page 61 says that its lists are
necessary subgroup lattices, without establishing realizability of every
listed shape. Our example contradicts necessity, so that qualification
does not remove the counterexample. We do not claim the geometric
intersection property. Nor is this rank-three example locally doubly
transitive: in the local action on a panel of size thirteen, the
stabilizer of one point has order four and cannot be transitive
on the other twelve points.

\section{A separate rank-two consequence}
The preceding example should be distinguished from the locally doubly
transitive geometries in \cite{scalars}.

\begin{corollary}\label{ranktwo}
The uniqueness assertion for composite $q-1$ in
\cite[Theorem 6.1, printed p.61]{source} is false at $q=32768$.
\end{corollary}
\begin{proof}
Here $q-1=32767=7\cdot31\cdot151$ is not prime.
The source asserts that the unique residually weakly primitive and
locally doubly transitive pregeometry has rank two and parabolics
$D_{2(q-1)}$ and $q:(q-1)$, with common subgroup of order $q-1$.
Its two valencies are $2$ and $q$.

Apply \cite[Theorem 3]{scalars} to $F=\F_{32768}$ and $K=\F_8$.
For each two-dimensional $K$-subspace $T$ containing $K$, it gives
a rank-two geometry with parabolics $\Sz(8)$ and $Q_T$,
where $|Q_T|=3584$ and their intersection has order $448$.
It is residually weakly primitive and locally doubly transitive, with
valencies $65$ and $8$. These differ from $2$ and $32768$, even after
interchanging types. Thus even the underlying incidence geometry is
absent from the uniqueness assertion. The forty-one classes counted
in \cite[Corollary 5]{scalars} use that paper's stated equivalence
under automorphisms of the specified group; a count of forty-one bare
incidence isomorphism classes is not needed or asserted here.
\end{proof}

\section{Attribution and verification}
The binary algebra, fixed scalars, matrices and rank-three construction
were obtained in the project's earlier manuscript on compatible Suzuki
tori, and were separately reviewed there. We retain that proof, its
certificate, original program and review as provenance. The present
note reuses the exact witness; its intersection proof uses group orders
in place of the earlier enumeration of $2132$ elements. The short new
replay reads the fixed witness, verifies (3)--(6) and the arithmetic
used in Lemma~\ref{meets}, and never searches for $M$ or enumerates an
ambient group. Classical Suzuki structure is retained explicitly.

The exact subfield-intersection lemma and the scalar-subspace theorem
are acquired from \cite{subfields,scalars}; their complete proofs and
reviews accompany this note. The rank-two corollary is a direct use of
the latter theorem, not a separate discovery. The earlier higher-rank
extension is not a premise. No other source cases or future realization
questions are declared settled, and no complete corrected classification,
publication priority, new conjecture solution or full Lean proof is claimed.

\begin{thebibliography}{9}
\bibitem{source} D. Leemans, \emph{The rank 3 geometries of the simple
Suzuki groups $\Sz(q)$}, Note di Matematica 19(1) (1999), 43--63.
\bibitem{subfields} Human Compression project,
\emph{A missing boundary in the Suzuki subfield construction},
reviewed project manuscript, 2026, Lemma 1.
\bibitem{scalars} Human Compression project,
\emph{Scalar subspaces in rank-two Suzuki geometries},
reviewed project manuscript, 2026, Theorem 3 and Corollary 5.
\end{thebibliography}
\end{document}
